%
% spec-mission-lifecycle.tex
%
% Mission Lifecycle: A Z Specification
% A formal model of the mission-contract state machine implemented in
% internal/mission (mission.go, store.go, validate.go, delegation.go,
% correct.go) and, for the two delegation-refusal call sites Abandon's
% gate depends on, internal/hook (pretooluse_dispatch.go,
% subagent_start.go). Current as of commit 244b6c3 (PR #517), the most
% recent commit to change lifecycle-affecting logic in internal/mission
% -- see \S\ref{sec:currency} for why that commit itself required no
% amendment here. Re-baselined
% from the original 9a60b39 (PR #432) pin to absorb DES-072 (#478),
% DES-075 (#508), and DES-076 (#509) -- see \S\ref{sec:des072} and
% \S\ref{sec:opabandon}.
%
% Compile: pdflatex spec-mission-lifecycle.tex
%
% Must be .tex, not .md: fuzz and pdflatex read content regardless of
% extension, but this repo's own z-spec tooling (file discovery for
% /z-spec:test and friends) globs for .tex specifically -- a .md file
% here is silently invisible to it and ProB reports 0 states.
%

\documentclass[a4paper,10pt,fleqn]{article}
\usepackage[margin=1in]{geometry}
\usepackage{fuzz}
\usepackage{booktabs}
\usepackage{tabularx}
\usepackage[colorlinks=true,linkcolor=blue,citecolor=blue,urlcolor=blue]{hyperref}

\newcommand{\fpath}[1]{\texttt{#1}}
\newcommand{\field}[1]{\texttt{#1}}

\begin{document}

\title{Mission Lifecycle: A Z Specification}
\author{Mike S, for Punt Labs}
\date{August 2026 (re-baselined September 2026)}
\hypersetup{
  pdftitle={Mission Lifecycle: A Z Specification},
  pdfauthor={Punt Labs},
  pdfsubject={Z Specification},
  pdfcreator={fuzz/probcli}
}
\maketitle
\tableofcontents
\newpage

% ===================================================================
\section{Scope}
\label{sec:scope}
% ===================================================================

This specification models the mission-contract lifecycle as it is
implemented today in \fpath{internal/mission}: the status values a
contract can hold (\field{validStatuses}, mission.go:10--28), the nine
transitions named below --- create, delegate, delegate-aborted,
disclaim-delegation, submit-result, reflect, advance-round, close,
abandon --- and the invariants those transitions must jointly
preserve. \field{DelegateAborted} and \field{DisclaimDelegation} are
new in this revision (DES-076, \S\ref{sec:opabandon}); the other seven
carry over from the original baseline. It is a baseline, not a design
proposal: every predicate below is read off existing code, cited by
file and line, and the point of writing it down is to give a fixed
target that future changes to \fpath{internal/mission} must not
violate without a deliberate, reviewed amendment to this document ---
see the currency rule, \S\ref{sec:currency}.

\subsection{Modelling choice: one mission at a time}
\label{sec:modelling-one-mission}

The real \field{Store} manages many contracts concurrently, each
under its own per-mission \field{flock} (\field{withLock},
store.go:2455--2481; \field{withCreateLock}, store.go:2496--2599).
Every transition in this specification --- Create,
Delegate, DelegateAborted, DisclaimDelegation, SubmitResult, Reflect,
AdvanceRound, Close, Abandon --- is likewise a lock-scoped operation on
a single contract in the real code: \field{Store.Close}, for instance,
loads, mutates, and writes back exactly one contract file under
exactly one lock (store.go:956--1187). Nothing in the nine named
transitions reads or writes a \emph{second} mission's state. The one
exception is \field{Store.Create}'s cross-mission write-set conflict
scan (\field{checkWriteSetConflicts}, store.go:3139--3210), which is
explicitly out of scope --- see \S\ref{sec:nongoals}.

This licenses the central simplification of the model: the state
schema \field{Mission} (\S\ref{sec:state}) describes \emph{one}
contract's lifecycle-relevant state, not a store-wide partial
function from mission identifiers to contracts. Nothing here would
change if it were promoted to $\field{MISSIONID} \pfun \field{Mission}$
--- every transition still touches exactly one point in that
function --- so the identifier type is omitted for economy.

\subsection{Non-goals}
\label{sec:nongoals}

The following are deliberately excluded. Each is a real mechanism in
\fpath{internal/mission}, but none of them bears on the four
invariants this specification checks (result gate, abandon gate,
round budget, evaluator-hash pinning), and including them would
obscure rather than sharpen the model:

\begin{itemize}
  \item \textbf{Cross-mission write-set conflicts} ---
    \field{checkWriteSetConflicts} (store.go:3139--3210) and
    \field{findWriteSetConflicts} (conflict.go:70) compare a new
    contract's \field{write\_set} against every other \emph{open}
    mission's. This is a Create-time admission check over the whole
    store, not a property of one mission's lifecycle.
  \item \textbf{Role overlap} --- \field{checkRoleOverlap}
    (store.go:3010--3138) refuses a
    contract whose worker and evaluator share a team role binding. It
    depends on an external \field{RoleLister} and is opt-in (nil
    disables it); it is an identity-system concern, not a
    lifecycle-state concern.
  \item \textbf{Generic \field{Store.Update}} (store.go:873--955) lets
    a caller rewrite an arbitrary field of an open contract. It is not
    one of the nine named transitions and none of the CLI/MCP surfaces
    route ordinary lifecycle events through it.
  \item \textbf{Pipelines, \field{depends\_on}, delegation-inheritance
    templates, and preconditions} (\field{Contract} fields
    \field{Pipeline}, \field{DependsOn}, \field{Preconditions},
    \field{Delegations} --- mission.go:111--173) are advisory or
    admission-control concerns layered on top of the status machine;
    they do not change which statuses exist or which transitions
    between them are legal.
  \item \textbf{\field{Abandon}'s and \field{DisclaimDelegation}'s
    \field{reason} arguments} (store.go:1291--1296, 1517--1522) are
    validated for non-emptiness and absence of control characters, and
    are redacted (\field{NewPathRedactor}, store.go:1323--1326,
    1532--1535) before they touch disk --- all data-hygiene concerns
    with no bearing on any state invariant. The same class of check
    was widened repo-wide by PR \#517 (commit 244b6c3, ethos-i1z8/
    ethos-21v0): \field{Result}, \field{Reflection}, and
    \field{Correction} free-text fields are now refused when a YAML
    plain scalar carries a discarded \texttt{\#}-comment suffix
    (\fpath{validate.go}, \fpath{reflection.go},
    \fpath{correct.go} --- each field's own \field{Validate}). This is
    the same non-goal as the reason-argument bullet above, generalised:
    a refusal on the \emph{content} of a free-text field never changes
    which \field{Status} a mission can reach or which round it is on,
    so it is modelled here only as a side condition in prose, never as
    a schema component.
  \item \textbf{\field{ForceReleaseWriteSet}} (store.go:1673--1838,
    DES-052-adjacent, ethos-9x07) marks an open mission's
    \field{write\_set}/\field{extract\_into} as released from
    admission control (\field{WriteSetReleasedAt}) without touching
    \field{Status}, \field{CurrentRound}, delegations, or results ---
    ``\emph{not a terminal transition, and not a substitute for Abandon
    or Close}'' (store.go:1650--1652). It is, like
    \field{checkWriteSetConflicts}, an admission-control concern
    layered on the status machine, not a status-machine transition
    itself: no value in \field{Mission} (\S\ref{sec:state}) changes
    when it runs.
\end{itemize}

\subsection{DES-072 correction events: excluded, with cause}
\label{sec:des072}

DES-072 (DESIGN.md; \#478, commit-dated 2026-08-22) adds
\field{Store.Correct} (correct.go:187--228), a new operation on an
\emph{already-closed} mission: it records a factual correction, a
fabrication finding, or a decision annotation as a new \texttt{correct}
event on the mission's append-only event log. This is exactly the
kind of post-terminal write the mission brief for this re-baseline
asks to be either modelled or excluded by name, with a restated
\field{TerminalIsFinal} if modelled. It is excluded, for a reason the
implementation itself states as a design invariant, not merely as an
empirical observation about the current code:

\begin{quote}
``\emph{Correct never touches contract.yaml, results.yaml,
reflections.yaml, or any of the Mission schema's fields --- the
operation's entire effect is one appendEventLocked call with Event:
"correct". This is why Correct cannot falsify the TerminalIsFinal
theorem: there is no Mission value in scope for it to mutate.}''
(correct.go:163--171)
\end{quote}

Three facts, checked directly against \field{Store.Correct}
(correct.go:187--228), justify the exclusion:

\begin{enumerate}
  \item \textbf{Guard is the complement of every modelled operation's.}
    \field{Correct} refuses \field{status = stOpen}
    (correct.go:207--210: ``\emph{corrections apply only to closed
    missions}'') --- the opposite of every operation in
    \S\ref{sec:ops}, all of which \emph{require}
    $\field{status} = \field{stOpen}$. A correction is definitionally
    unreachable from any state this specification's operation set can
    still act on.
  \item \textbf{No write touches a \field{Mission} field.} The
    function body (correct.go:211--226) validates the correction,
    checks \field{staged.Mission} and \field{staged.Round} against the
    loaded contract, then calls exactly one
    \field{s.appendEventLocked} (correct.go:221--226) with
    \field{Event: EventCorrect}. \field{status}, \field{CurrentRound},
    \field{Budget.Rounds}, \field{Worker}, \field{Evaluator}, the
    delegations directory, and the results file are all read-only
    inputs to a validation check --- none is ever written.
  \item \textbf{The event log is not part of the \field{Mission}
    state schema.} \S\ref{sec:state}'s \field{Mission} schema already
    omits the event log entirely (it is an audit trail, not lifecycle
    state --- the same omission that lets \field{Leader}, per
    \S\ref{sec:basic}, be left out of the schema even though it rides
    on every event). A \texttt{correct} event appended to a log this
    specification does not model cannot change a reachable
    \field{Mission} value, by construction.
\end{enumerate}

Given all three, \field{TerminalIsFinal} (\S\ref{sec:theorems}) needs
no restatement: there is no new value \field{status} can take
(\field{Correct} never assigns one), and no additional way for
\field{status}$'$, \field{currentRound}$'$,
\field{blockingDelegationCount}$'$, or \field{resultRounds}$'$ to
differ from their pre-state once a mission is terminal. \field{Correct} is a real, callable,
CLI/MCP-exposed operation --- and a load-bearing fix for the
fabricated-result incident that motivated it (ethos-11fy, ethos-ecpv)
--- but it operates entirely outside the state space this
specification checks, which is exactly the intended shape: an audit
annotation should never be able to reopen, or even perturb, a
finished mission's status machine.

\subsection{Currency}
\label{sec:currency}

This specification is a baseline, and baselines drift silently unless
something calls them back into review. Nothing in CI or in a git hook
enforces this mechanically today --- the rule below is a review-
checklist convention, the same status as every other entry in
\fpath{docs/development.md}'s Standards Checklist, not a gate that
blocks a merge on its own. The rule: \textbf{a pull
request that changes lifecycle-affecting logic in
\fpath{internal/mission} --- a new \field{Status} or
\field{Recommendation} value, a changed guard on \field{Create},
\field{Update}, \field{AppendResult}, \field{AppendReflection},
\field{AdvanceRound}, \field{Close}, \field{Abandon}, or
\field{DisclaimDelegation}, or a new operation with a status-machine
effect --- must either amend this document in the same PR, or state
explicitly in the PR description why the change does not affect any
schema, operation, or theorem here.} The companion rule lives in
\fpath{docs/development.md}'s Standards Checklist, so it surfaces
at the point a worker is deciding what a PR needs before it can merge,
not only here where it is easy to forget mid-implementation.

This re-baseline is itself evidence for the rule's necessity: PR
\#509 (DES-076, commit 5963093) changed \field{Abandon}'s gate
predicate --- exactly the shape of change the rule now requires
flagging --- and this document was not amended until three PRs and
twelve days later, on the discovery that the on-disk pin (9a60b39,
PR \#432) was three lifecycle-affecting PRs stale
(\#478 DES-072, \#508 DES-075, \#509 DES-076). PR \#517 (commit
244b6c3), by contrast, is the rule's other branch working correctly:
it widened free-text validation (\S\ref{sec:nongoals}) without
touching any \field{Status}, guard, or transition this specification
tracks, so the currency rule's second clause applies and no schema
amendment was owed --- recorded here, per the rule's own requirement,
rather than left to be re-discovered by the next reader who diffs the
pinned commit against \fpath{internal/mission}'s history.

% ===================================================================
\section{Basic Types}
\label{sec:basic}
% ===================================================================

Two given sets carry no internal structure relevant to the lifecycle:
identities (leader, worker, evaluator handles all draw from the same
set in the real code --- \field{Contract.Leader},
\field{Contract.Worker}, \field{Contract.Evaluator.Handle} are all
plain \field{string} handles, mission.go:52--54) and the evaluator's
pinned content hash (a SHA-256 hex digest, \field{Evaluator.Hash},
mission.go:223).

\begin{zed}
[HANDLE, HASH]
\end{zed}

\field{Leader} is carried on every event as the acting identity
(store.go:1021, 1602, 1809 (via reflect/result authors), 2944) but
plays no role in any lifecycle invariant, so it is omitted from the
state schema for economy --- see \S\ref{sec:state}.

% ===================================================================
\section{Free Types}
\label{sec:freetypes}
% ===================================================================

\subsection{Status}

The five values of \field{validStatuses} (mission.go:10--28,
validate.go:28--34), prefixed \texttt{st} to keep the constructors
distinct from any B keyword ProB might reserve:

\begin{zed}
Status ::= stOpen | stClosed | stFailed | stEscalated | stAbandoned
\end{zed}

\field{stClosed}, \field{stFailed}, and \field{stEscalated} are the
three statuses \field{Store.Close} can produce
(store.go:963--964: ``\emph{must be closed, failed, or escalated}'').
\field{stAbandoned} is reachable only through \field{Store.Abandon}
(store.go:1290--1487) and is explicitly excluded from \field{Close}'s
own range (store.go:957--965, comment: ``\emph{StatusAbandoned is
excluded here on purpose}'').

\begin{axdef}
CloseableStatus : \power Status
\where
CloseableStatus = \{ stClosed, stFailed, stEscalated \}
\end{axdef}

\subsection{Recommendation}

The five values of \field{validRecommendations}
(reflection.go:89--103), governing whether \field{AdvanceRound} may
fire (store.go:2916--2927, \field{IsTerminalRecommendation}):

\begin{zed}
Recommendation ::= recContinue | recAdvance | recPivot
                  | recStop | recEscalate
\end{zed}

\begin{axdef}
TerminalRecommendation : \power Recommendation
\where
TerminalRecommendation = \{ recStop, recEscalate \}
\end{axdef}

\field{recContinue} and \field{recAdvance} are synonyms in the real
code (reflection.go:85--88: ``\emph{continue and advance are synonyms
for `advance permitted, same approach'}''); both are excluded from
\field{TerminalRecommendation}, matching
\field{IsTerminalRecommendation}'s two-value check exactly
(reflection.go:111--113).

\subsection{ZBOOL}

fuzz has no built-in boolean, and \field{BOOL}/\field{true}/
\field{false} collide with B keywords under ProB's translation:

\begin{zed}
ZBOOL ::= ztrue | zfalse
\end{zed}

This carries one piece of state-external configuration: whether the
\field{Store} the mission is operating under was constructed with a
non-empty \field{repoRoot} (store.go:218--231, 1374--1380). This is a
property of the \emph{Store instance}, not of the contract on disk
--- the same contract can be abandoned successfully by one caller and
refused by another, depending only on how each caller's \field{Store}
was built. \S\ref{sec:validation} depends on this distinction being
explicit.

% ===================================================================
\section{Global Constants}
\label{sec:constants}
% ===================================================================

The real \field{minRounds}/\field{maxRounds} bound is $[1, 10]$
(validate.go:228--231, 243--253, rules for \field{budget.rounds} and
\field{current\_round} respectively). This specification uses a
smaller bound purely to keep the state space finite for ProB (per the
z-spec ProB-compatibility guidance on bounding numeric variables) ---
the shape of every invariant below is independent of the concrete
value.

\begin{axdef}
maxRounds : \nat
\where
maxRounds = 3
\end{axdef}

The real code imposes no maximum on how many delegations a mission may
accumulate: \field{countBlockingDelegations} (store.go:1959--1999),
the function \field{blockingDelegationCount} models, returns an
unbounded count, as does its sibling \field{countDelegations}
(store.go:1856--1875 --- the un-excluded total, used only by
\field{ForceReleaseWriteSet}'s informational staleness snapshot,
\S\ref{sec:nongoals}). This specification bounds
\field{blockingDelegationCount} only so ProB can enumerate the state
space; the invariant in \S\ref{sec:state} does not depend on the
bound's value, only on the fact that \emph{some} finite count is
tracked.

\begin{axdef}
maxDelegations : \nat
\where
maxDelegations = 2
\end{axdef}

% ===================================================================
\section{State}
\label{sec:state}
% ===================================================================

\begin{schema}{Mission}
status : Status \\
currentRound : \nat \\
budgetRounds : \nat \\
worker, evaluatorHandle : HANDLE \\
evaluatorHash : HASH \\
blockingDelegationCount : \nat \\
resultRounds : \finset \nat \\
reflected : \nat \pfun Recommendation \\
repoRootSet : ZBOOL
\where
budgetRounds \in 1 \upto maxRounds \\
currentRound \in 1 \upto budgetRounds \\
blockingDelegationCount \leq maxDelegations \\
resultRounds \subseteq 1 \upto budgetRounds \\
\dom reflected \subseteq 1 \upto budgetRounds \\
status = stAbandoned \implies
  (blockingDelegationCount = 0 \land resultRounds = \emptyset)
\end{schema}

\field{blockingDelegationCount} models \field{countBlockingDelegations}'s
return value (store.go:1959--1999) --- the count \field{Abandon}'s own
gate reads --- \emph{not} the raw number of delegation records on disk
(that is \field{countDelegations}, store.go:1856--1875, excluded per
\S\ref{sec:nongoals} because nothing but the non-goal
\field{ForceReleaseWriteSet} ever reads it). This is a deliberate
renaming from the original baseline's plain \field{delegationCount}:
DES-076 makes ``how many delegation records exist'' and ``how many
still block Abandon'' different questions
(store.go:1882--1886: ``\emph{`how many entries' and `how many still
block' are no longer the same question once either exclusion
applies}''), and only the second one bears on any invariant this
specification checks.

Six predicates, grouped into three invariants:

\begin{itemize}
  \item \field{budgetRounds} $\in 1 \upto \field{maxRounds}$ and
    \field{currentRound} $\in 1 \upto \field{budgetRounds}$ together
    are validate.go rule~12 (\field{budget.rounds} $\in [1,10]$,
    validate.go:228--231) and rule~14 (\field{current\_round} $\in
    [1,\field{budget.rounds}]$, validate.go:243--253) --- \textbf{round
    number never exceeds the budget}.
  \item \field{blockingDelegationCount} $\leq \field{maxDelegations}$
    and the \field{resultRounds}/\field{reflected} domain bounds are
    ProB animatability conditions with no code counterpart --- they
    only keep the state space finite.
  \item The final line is the \textbf{abandon invariant}: a mission
    in state \field{stAbandoned} has zero \emph{blocking} delegations
    and no results for any round. This is the property store.go:
    1234--1289's doc comment states as Abandon's premise (``\emph{was
    there ever any work to lose?}'', store.go:1248--1250) and
    store.go:1422--1441 checks for \emph{every} round, not only the
    current one. Unlike the pre-DES-076 baseline, this does
    \emph{not} mean an abandoned mission's delegations directory is
    necessarily empty --- an aborted or disclaimed delegation record
    can remain on disk (\S\ref{sec:opabandon}); the invariant is about
    what \emph{blocks}, not what \emph{exists}.
\end{itemize}

Two invariants from the brief are \emph{not} state predicates and so
do not appear in this schema:

\begin{itemize}
  \item \textbf{Close requires a result for the current round.} This
    is a precondition on the \emph{transition into} a terminal
    status, not a property every reachable state must satisfy in
    isolation (an open mission with no result yet is perfectly
    valid). It appears as a guard on \field{Close}
    (\S\ref{sec:opclose}).
  \item \textbf{The evaluator hash never changes.} This is a
    property of a \emph{trajectory} through the state space (every
    operation's after-state agrees with its before-state on
    \field{evaluatorHash}), not a property of any single state. Z has
    no first-class temporal-invariant construct; it is stated as an
    explicit frame conjunct, \field{evaluatorHash}$'$ =
    \field{evaluatorHash}, in every operation below, and recorded as
    a theorem in \S\ref{sec:theorems}.
\end{itemize}

\field{worker} $\neq$ \field{evaluatorHandle} is likewise absent from
this schema. \field{checkSelfVerification} (store.go:2972--3009)
enforces it, but only at \field{Store.Create} --- it is not one of
\field{Contract.Validate}'s numbered rules (validate.go:41--97), so a
hand-edited on-disk contract with \field{worker} = \field{evaluator}
would load without error. Encoding it as a standing schema invariant
would overstate what the real code guarantees. It appears instead as
a guard on \field{Init} only (\S\ref{sec:init}); every other
operation below leaves \field{worker} and \field{evaluatorHandle}
unchanged, so the property, once established at creation, survives
every transition this specification models --- but that is a fact
about which operations exist, not a checked invariant.

% ===================================================================
\section{Initialization}
\label{sec:init}
% ===================================================================

\field{Init} corresponds to \field{Store.Create}
(store.go:606--728) as prepared by \field{ApplyServerFields}
(store.go:323--368): a freshly created mission is always
\field{stOpen}, at round~1, with zero blocking delegations and zero
results.

\begin{schema}{Init}
Mission'
\where
status' = stOpen \\
currentRound' = 1 \\
worker' \neq evaluatorHandle' \\
blockingDelegationCount' = 0 \\
resultRounds' = \emptyset \\
reflected' = \emptyset
\end{schema}

\field{budgetRounds}$'$, \field{evaluatorHash}$'$, and
\field{repoRootSet}$'$ are left unconstrained by \field{Init} beyond
\field{Mission}'s own type and range predicates: the caller chooses
the round budget (mission.go:241--243), the hash is computed from
live identity data before \field{Create} is even called
(store.go:314--318), and \field{repoRootSet} reflects how the caller
built their \field{Store}, not a choice \field{Create} itself makes.

% ===================================================================
\section{Operations}
\label{sec:ops}
% ===================================================================

Every operation below requires \field{status = stOpen} as a
precondition, with the sole exception of \field{Idle}
(\S\ref{sec:idle}), whose entire purpose is to be enabled exactly
where the others are not. This is not incidental repetition: it is
the single fact from which \textbf{terminal states are final} follows
--- \field{Store.Update}, \field{AppendResult},
\field{AppendReflection}, \field{AdvanceRound}, \field{Close},
\field{Abandon}, and \field{DisclaimDelegation} all begin with the
same shape of check (``\emph{mission \%q is in terminal state \%q}'',
store.go:985--987, 1341--1346, 1542--1547, 2891--2893, 3360--3362,
2727--2729). \field{Delegate} and \field{DelegateAborted} carry the
same guard as a modelling choice (\S\ref{sec:opdelegate}). Once a
mission leaves \field{stOpen}, nothing in this specification --- and
nothing in the real \field{Store} --- can move it again.

\subsection{Delegate}
\label{sec:opdelegate}

Models a worker spawn recorded via \field{WriteDelegationSkeleton}
(delegation.go:346--396), which always stamps the new record
\field{Verdict: DelegationVerdictOpen} (delegation.go:391) --- every
successful spawn is, at the moment it is recorded, a \emph{blocking}
delegation; nothing distinguishes it from one that will later close
\field{pass}/\field{fail}/\field{escalate}. The real function has no
status check of its own --- it is called by the leader's orchestration,
not gated by the store --- so the \field{status = stOpen} guard here is
a modelling choice: it is not sound, operationally, to spawn a worker
against a mission that is already closed, and every real call site
only ever delegates against an open mission.

\begin{schema}{Delegate}
\Delta Mission
\where
status = stOpen \\
blockingDelegationCount < maxDelegations \\
status' = status \\
currentRound' = currentRound \\
budgetRounds' = budgetRounds \\
worker' = worker \\
evaluatorHandle' = evaluatorHandle \\
evaluatorHash' = evaluatorHash \\
blockingDelegationCount' = blockingDelegationCount + 1 \\
resultRounds' = resultRounds \\
reflected' = reflected \\
repoRootSet' = repoRootSet
\end{schema}

\subsection{DelegateAborted}
\label{sec:opdelegateaborted}

New in this revision. Models a delegation that is spawned and then
refused \emph{before its worker ever runs}, collapsed into one atomic
step: two real events in sequence, \field{WriteDelegationSkeleton}
(delegation.go:346, \field{Verdict: DelegationVerdictOpen}) followed
immediately, within the same synchronous hook invocation, by
\field{CloseDelegationSkeleton} (delegation.go:649--726) stamping
\field{Verdict: DelegationVerdictAborted}. Exactly two call sites do
this, and \field{TestAbortedVerdictWriteSites\_MatchKnownThree}
(store.go:1907--1914, referring to
\fpath{aborted\_writer\_drift\_test.go}) enforces that the set stays
exactly three (the third, \field{closeDelegationSkeletons}'
escalate-mapping inside \field{Close}, cannot fire here --- see below):

\begin{itemize}
  \item the \field{max\_delegation\_depth} refusal,
    \field{closeDelegationAborted}
    (internal/hook/pretooluse\_dispatch.go:991--1014);
  \item the content-hash verifier-gate refusal,
    \field{closeSkeletonOnHashRefusal}
    (internal/hook/subagent\_start.go:787--833).
\end{itemize}

\field{countBlockingDelegations} (store.go:1959--1999) skips any
record with \field{Verdict == DelegationVerdictAborted}
(store.go:1993--1995) unconditionally --- independent of
\field{BoundVia} or any disclaim (store.go:1888--1889: ``\emph{does
NOT depend on BoundVia or a disclaim}''). Both call sites above fire
strictly before the delegated worker process starts
(store.go:1890--1905), so every delegation this schema represents is,
from Abandon's point of view, indistinguishable from one that was
never spawned at all.

\begin{schema}{DelegateAborted}
\Xi Mission
\where
status = stOpen
\end{schema}

$\Xi \field{Mission}$ --- no field changes, not even
\field{blockingDelegationCount} --- is not a simplification of this
operation's effect; it \emph{is} the effect. DES-076's guarantee for
the aborted case is precisely that this event is invisible to every
field the \field{Mission} schema tracks, so the honest schema is a
self-loop, structurally the mirror image of \field{Idle}
(\S\ref{sec:idle}): \field{Idle} is enabled exactly where
$\field{status} \neq \field{stOpen}$, \field{DelegateAborted} exactly
where $\field{status} = \field{stOpen}$, and neither ever moves the
state.

\subsection{DisclaimDelegation}
\label{sec:opdisclaim}

New in this revision. Models \field{Store.DisclaimDelegation}
(store.go:1516--1633), which excludes one operator-named, already-
blocking delegation from \field{blockingDelegationCount} without
weakening the count for any other delegation. The real per-delegation
eligibility check, \field{DisclaimDelegationRecord}
(delegation.go:726--790), is mechanical and three-part
(delegation.go:707--720): \field{BoundVia} must equal
\field{BoundViaSidecarDispatch} exactly; \field{Verdict} must not be
\field{DelegationVerdictOpen} (the delegation must already be closed);
and the record must not already be disclaimed. None of the three is
representable at this specification's granularity, which --- per the
\S\ref{sec:modelling-one-mission} economy argument already applied to
the contract as a whole --- tracks delegations only as a count, not as
individual records with provenance and verdict fields. The guard below
is the coarsest sound abstraction of that three-part check: a
disclaim needs \emph{some} still-blocking delegation to name, exactly
as \field{DisclaimDelegationRecord} needs the one it is given to exist
and be eligible; the specific eligibility distinctions
(docs/mission-abandon.md:92--110) are a property of \emph{which}
delegation is picked, not of whether the operation can fire at all,
and so are omitted here the same way \field{Abandon}'s \field{reason}
argument is (\S\ref{sec:nongoals}).

\begin{schema}{DisclaimDelegation}
\Delta Mission
\where
status = stOpen \\
blockingDelegationCount > 0 \\
status' = status \\
currentRound' = currentRound \\
budgetRounds' = budgetRounds \\
worker' = worker \\
evaluatorHandle' = evaluatorHandle \\
evaluatorHash' = evaluatorHash \\
blockingDelegationCount' = blockingDelegationCount - 1 \\
resultRounds' = resultRounds \\
reflected' = reflected \\
repoRootSet' = repoRootSet
\end{schema}

The status guard here is not merely a modelling choice, unlike
\field{Delegate}'s: the real \field{Store.DisclaimDelegation} checks
$\field{c.Status} \neq \field{StatusOpen}$ explicitly and refuses
(store.go:1542--1547, ``\emph{disclaim only applies to an open mission
awaiting abandon}'').

\subsection{SubmitResult}

Models \field{Store.AppendResult} (store.go:3334--3450): a result is
accepted only on an open mission, only for the current round, and
never twice for the same round (store.go:3360--3362, the append-only
check at store.go:3393--3400).

\begin{schema}{SubmitResult}
\Delta Mission \\
round? : \nat
\where
status = stOpen \\
round? = currentRound \\
round? \notin resultRounds \\
status' = status \\
currentRound' = currentRound \\
budgetRounds' = budgetRounds \\
worker' = worker \\
evaluatorHandle' = evaluatorHandle \\
evaluatorHash' = evaluatorHash \\
blockingDelegationCount' = blockingDelegationCount \\
resultRounds' = resultRounds \cup \{ round? \} \\
reflected' = reflected \\
repoRootSet' = repoRootSet
\end{schema}

\subsection{Reflect}

Models \field{Store.AppendReflection} (store.go:2703--2807): same
shape as \field{SubmitResult} --- open mission, current round only,
append-only (store.go:2727--2729, 2749--2756).

\begin{schema}{Reflect}
\Delta Mission \\
round? : \nat \\
rec? : Recommendation
\where
status = stOpen \\
round? = currentRound \\
round? \notin \dom reflected \\
status' = status \\
currentRound' = currentRound \\
budgetRounds' = budgetRounds \\
worker' = worker \\
evaluatorHandle' = evaluatorHandle \\
evaluatorHash' = evaluatorHash \\
blockingDelegationCount' = blockingDelegationCount \\
resultRounds' = resultRounds \\
reflected' = reflected \cup \{ round? \mapsto rec? \} \\
repoRootSet' = repoRootSet
\end{schema}

\subsection{AdvanceRound}

Models \field{Store.AdvanceRound} (store.go:2878--2963): the budget
check comes first in the real code (``\emph{the operator sees the
right diagnostic}'', store.go:2894--2898), then the reflection-exists
check, then the terminal-recommendation refusal
(store.go:2916--2927).

\begin{schema}{AdvanceRound}
\Delta Mission
\where
status = stOpen \\
currentRound < budgetRounds \\
currentRound \in \dom reflected \\
reflected(currentRound) \notin TerminalRecommendation \\
status' = status \\
currentRound' = currentRound + 1 \\
budgetRounds' = budgetRounds \\
worker' = worker \\
evaluatorHandle' = evaluatorHandle \\
evaluatorHash' = evaluatorHash \\
blockingDelegationCount' = blockingDelegationCount \\
resultRounds' = resultRounds \\
reflected' = reflected \\
repoRootSet' = repoRootSet
\end{schema}

\subsection{Close}
\label{sec:opclose}

Models \field{Store.Close} (store.go:956--1187). The result gate,
\field{checkResultGateLocked} (store.go:3516--3533), is the
\textbf{close-requires-a-result} invariant from the brief: it refuses
unconditionally unless a result exists for exactly the mission's
current round (store.go:3522: ``\emph{Round == c.CurrentRound}'').

\begin{schema}{Close}
\Delta Mission \\
newStatus? : Status
\where
status = stOpen \\
newStatus? \in CloseableStatus \\
currentRound \in resultRounds \\
status' = newStatus? \\
currentRound' = currentRound \\
budgetRounds' = budgetRounds \\
worker' = worker \\
evaluatorHandle' = evaluatorHandle \\
evaluatorHash' = evaluatorHash \\
blockingDelegationCount' = blockingDelegationCount \\
resultRounds' = resultRounds \\
reflected' = reflected \\
repoRootSet' = repoRootSet
\end{schema}

\subsection{Abandon}
\label{sec:opabandon}

Models \field{Store.Abandon} (store.go:1290--1487; doc comment
store.go:1234--1289) as it stands \emph{today} --- after both the
repoRoot fail-open fix tested by
\field{TestStore\_Abandon\_RefusesWhenRepoRootEmpty}
(abandon\_test.go:151--179, \S\ref{sec:validation}) \emph{and} DES-076
(\#509, commit 5963093), which replaced the original ``zero
delegation records, full stop'' gate with a narrower, mechanically
gated count. Four conjuncts, four gates in the real code:

\begin{itemize}
  \item \field{repoRootSet = ztrue} --- store.go:1374--1380: an empty
    \field{repoRoot} \textbf{refuses}, unconditionally, before any of
    the other gates is even evaluated.
  \item \field{status = stOpen} --- store.go:1341--1346: an
    already-terminal mission refuses re-abandonment for the same
    reason \field{Close} refuses re-closing.
  \item \field{blockingDelegationCount = 0} --- Gate~1,
    store.go:1347--1421, backed by \field{countBlockingDelegations}
    (store.go:1959--1999), which excludes exactly the two DES-076
    carve-outs modelled as their own operations above
    (\S\ref{sec:opdelegateaborted}, \S\ref{sec:opdisclaim}): a
    delegation with \field{Verdict == DelegationVerdictAborted}
    (store.go:1993--1995) and one with a non-empty
    \field{DisclaimedAt} (store.go:1990--1992). See also
    docs/mission-abandon.md:39--59 for the gate's narrative form and
    the two carve-outs' rationale.
  \item \field{resultRounds = \emptyset} --- Gate~2, store.go:1422
    --1441, checked for \emph{every} round on record, not only the
    current one.
\end{itemize}

\begin{schema}{Abandon}
\Delta Mission
\where
status = stOpen \\
repoRootSet = ztrue \\
blockingDelegationCount = 0 \\
resultRounds = \emptyset \\
status' = stAbandoned \\
currentRound' = currentRound \\
budgetRounds' = budgetRounds \\
worker' = worker \\
evaluatorHandle' = evaluatorHandle \\
evaluatorHash' = evaluatorHash \\
blockingDelegationCount' = blockingDelegationCount \\
resultRounds' = resultRounds \\
reflected' = reflected \\
repoRootSet' = repoRootSet
\end{schema}

The \field{Mission} state invariant already requires
$\field{status}' = \field{stAbandoned} \implies
(\field{blockingDelegationCount}' = 0 \land \field{resultRounds}' =
\emptyset)$; \field{Abandon}'s own guard is not logically necessary
given that invariant and the frame equations above --- it is written
out anyway, because the guard is the design step (precondition
calculation is what \field{Store.Abandon}'s author had to get right),
and because \S\ref{sec:validation} needs a second operation, without
the guard's \field{repoRootSet} conjunct in the right place, to fail
the same invariant.

\emph{On what changed since the original baseline.} The pre-DES-076
gate was $\field{delegationCount} = 0$, where \field{delegationCount}
counted \emph{every} record under \fpath{delegations/} with no
exclusion at all --- so any spawned worker, whatever its eventual
verdict, blocked Abandon forever, and the only recourse for a mission
whose sole delegation was refused pre-run (or later proven to be a
dispatch-sidecar misattribution) was none: \field{Close} correctly
refuses without a result, and Abandon refused too. DES-076's fix is
visible in this specification as a difference in \emph{which count}
the same-shaped guard checks, plus the two new operations
(\S\ref{sec:opdelegateaborted}, \S\ref{sec:opdisclaim}) that can bring
that count back to zero after a \field{Delegate} without ever
producing a result --- exactly the two carve-outs
docs/mission-abandon.md:39--59 names, and no others: a delegation that
ran to a normal \field{pass}/\field{fail}/\field{escalate} close still
counts as blocking, matching store.go:1263--1268's ``\emph{a
non-aborted, non-disclaimed closed delegation ... refuses the
transition and points the caller at Close instead}''.

\subsection{Idle}
\label{sec:idle}

Not one of the transitions named in \S\ref{sec:scope}, and not a real
operation on \field{Store} at all. Every operation above except
\field{DelegateAborted} guards on $\field{status} = \field{stOpen}$
(\field{TerminalIsFinal}, \S\ref{sec:theorems}), so a state space
built only from those transitions has no outgoing transition once
\field{status} leaves \field{stOpen} --- ProB reports that correctly,
but as a \emph{deadlock}, which is a different finding from an
invariant violation and would mask the one \S\ref{sec:validation} is
looking for. \field{Idle} is the standard fix: a state-preserving
self-loop enabled exactly on terminal states, so a finished mission is
modelled as ``nothing more happens,'' not ``the machine got stuck.''

\begin{schema}{Idle}
\Xi Mission
\where
status \neq stOpen
\end{schema}

% ===================================================================
\section{Theorems}
\label{sec:theorems}
% ===================================================================

fuzz type-checks schemas; it does not discharge proof obligations.
The following are claims about the operation set
$\{\field{Init}, \field{Delegate}, \field{DelegateAborted},
\field{DisclaimDelegation}, \field{SubmitResult}, \field{Reflect},
\field{AdvanceRound}, \field{Close}, \field{Abandon}\}$, each backed
by inspection of the schemas above; two (\field{TerminalIsFinal},
\field{AbandonSafety}) are additionally exercised by ProB animation
(\S\ref{sec:validation}). \field{EvaluatorHashFrozen} is a syntactic
guarantee of every operation's frame conjunct --- no reachable
violation is even expressible, so animation would confirm nothing
non-trivial. \field{CloseRequiresResult}'s justification is a
file:line correspondence to the real Go source, which ProB cannot
check --- it never reads Go code. None of the four claims are
discharged by proof:

\begin{center}
\begin{tabularx}{\textwidth}{@{}lX@{}}
\toprule
\textbf{Theorem} & \textbf{Justification} \\
\midrule
\field{EvaluatorHashFrozen} &
  Every operation's predicate includes
  $\field{evaluatorHash}' = \field{evaluatorHash}$ (or, for
  \field{DelegateAborted} and \field{Idle}, leaves it untouched via
  $\Xi \field{Mission}$). No operation in the set assigns any other
  value, so \field{evaluatorHash} is constant on every reachable
  trajectory from \field{Init}. Matches store.go:314--318's comment
  that the hash is ``\emph{the trust anchor}'' pinned once at
  launch. \\
\field{TerminalIsFinal} &
  \field{Delegate}, \field{DelegateAborted}, \field{DisclaimDelegation},
  \field{SubmitResult}, \field{Reflect}, \field{AdvanceRound},
  \field{Close}, and \field{Abandon} all guard on $\field{status} =
  \field{stOpen}$. Once $\field{status}' \neq \field{stOpen}$ (via
  \field{Close} or \field{Abandon}), no operation in the set is
  enabled again --- \field{status} is absorbing at every value in
  \field{CloseableStatus} $\cup \{ \field{stAbandoned} \}$. DES-072's
  \field{Correct} does not extend this set; \S\ref{sec:des072}
  establishes why it cannot falsify this claim without being modelled
  at all. \\
\field{CloseRequiresResult} &
  \field{Close}'s guard $\field{currentRound} \in
  \field{resultRounds}$ is exactly the calculated precondition of
  \field{checkResultGateLocked} (store.go:3516--3533): the gate holds
  iff a result was recorded for the round the mission is closing
  \emph{from}. \\
\field{AbandonSafety} &
  The \field{Mission} invariant's final conjunct
  ($\field{status} = \field{stAbandoned} \implies
  \field{blockingDelegationCount} = 0 \land \field{resultRounds} =
  \emptyset$) holds in \field{Init}'s image and is preserved by every
  operation in the correct set: \field{Abandon} because its guard
  entails the consequent directly; \field{DelegateAborted} because it
  leaves \field{blockingDelegationCount} unchanged by definition
  ($\Xi \field{Mission}$); \field{DisclaimDelegation} because it only
  ever decreases \field{blockingDelegationCount}, never increases it;
  and every other operation because it either leaves \field{status} at
  \field{stOpen} or does not touch \field{status} at \field{stAbandoned}
  at all. \\
\bottomrule
\end{tabularx}
\end{center}

% ===================================================================
\section{Validation: the abandon fail-open regression}
\label{sec:validation}
% ===================================================================

This specification also demonstrates a concrete case where it would
have caught a real regression: the fail-open bug fixed in PR \#432
(commit 9a60b39) and regression-tested by
\field{TestStore\_Abandon\_RefusesWhenRepoRootEmpty}
(abandon\_test.go:151--179): the historical \field{Store.Abandon}
guarded its delegation and result checks with \texttt{if
s.repoRoot != ""}, so when \field{repoRoot} was empty, both checks
were \emph{skipped}, not refused --- a mission with real, recoverable
work on record could be abandoned cleanly, discarding it silently.
This demonstration predates DES-076 by design: the bug it exhibits is
about \field{repoRootSet}'s conjunct short-circuiting the \emph{whole}
delegation gate, independent of which count that gate happens to
check, so the DES-076 renaming (\field{delegationCount} $\to$
\field{blockingDelegationCount}, \S\ref{sec:state}) carries through
into this section as a variable rename only --- the technique, the
counter-example, and the errata below are otherwise unchanged from
the original baseline.

\subsection{AbandonFailOpen}

Not part of the current implementation, and not a schema in this
document's own checked namespace either --- for reasons the next
subsection explains. Its guard replaces \field{Abandon}'s independent
conjuncts with the historical code's single implication, reproducing
the \texttt{if repoRoot != ""} short-circuit exactly:

\begin{verbatim}
AbandonFailOpen
  Delta Mission
where
  status = stOpen
  (repoRootSet = ztrue ==>
    (blockingDelegationCount = 0 and resultRounds = {}))
  status' = stAbandoned
  currentRound' = currentRound
  budgetRounds' = budgetRounds
  worker' = worker
  evaluatorHandle' = evaluatorHandle
  evaluatorHash' = evaluatorHash
  blockingDelegationCount' = blockingDelegationCount
  resultRounds' = resultRounds
  reflected' = reflected
  repoRootSet' = repoRootSet
\end{verbatim}

When \field{repoRootSet = zfalse}, the guard's implication is
vacuously true regardless of \field{blockingDelegationCount} or
\field{resultRounds} --- exactly the fail-open condition djb's probe
found. The frame equations carry \field{blockingDelegationCount} and
\field{resultRounds} through \emph{unchanged}, matching the real bug:
the historical code did not clear or lose the delegation records, it
simply stopped checking them before transitioning \field{status} to
\field{stAbandoned}.

\subsection{Errata: why this schema is not checked against
\field{Mission}}

An earlier revision of this document defined \field{AbandonFailOpen}
as a live schema typed against $\Delta \field{Mission}$, alongside
\field{Abandon} and \field{Idle}, and claimed that model-checking the
combined operation set would surface the invariant violation directly.
It does not, and the claim was wrong. Model-checking
$\{\field{Init}, \field{Delegate}, \field{DelegateAborted},
\field{DisclaimDelegation}, \field{SubmitResult},
\field{Reflect}, \field{AdvanceRound}, \field{Close}, \field{Abandon},
\field{AbandonFailOpen}, \field{Idle}\}$ against \field{Mission}
explored the state space reported in the Tool run table
(\S\ref{sec:probrun}) and reported \textbf{zero} violations --- with
the buggy operation present and enabled.

The reason is the schema calculus, not a tooling gap. $\Delta
\field{Mission}$ abbreviates $\field{Mission} \land \field{Mission}'$,
and \field{Mission}$'$ is \field{Mission} with every component ---
\emph{and its own \texttt{where}-clause predicate} --- decorated with
a prime. \field{Mission}'s predicate includes the abandon invariant,
so \emph{every} schema written as $\Delta \field{Mission} \land P$
already carries, as a silent extra conjunct,
$\field{status}' = \field{stAbandoned} \implies
(\field{blockingDelegationCount}' = 0 \land \field{resultRounds}' =
\emptyset)$
--- regardless of what $P$ says. Combined with
\field{AbandonFailOpen}'s own frame equations
($\field{blockingDelegationCount}' = \field{blockingDelegationCount}$,
$\field{resultRounds}' = \field{resultRounds}$), this makes the whole
schema unsatisfiable at exactly the states that would exhibit the
bug: there is no valuation of the primed variables that satisfies both
the frame equations and the silent invariant conjunct when
$\field{blockingDelegationCount} > 0$. The operation is not
\emph{disabled} by a guard that checks for this --- it has no valid
after-state at all, by pure predicate logic, so ProB correctly finds
nothing wrong, because there genuinely is nothing there to find.
Writing a bad guard inside $\Delta \field{Mission}$ cannot, by
construction, violate an invariant already living in \field{Mission}'s
own predicate --- that is a structural property of Z's schema
calculus, not a claim about the real \field{Store.Abandon}.

This is precisely why \field{Contract.Validate} not carrying the
abandon invariant (\S\ref{sec:state}) turns out to matter for more
than fidelity to the real rule list: had this specification instead
baked \field{AbandonSafety} into a Contract-level type, \emph{no}
buggy guard could ever be exhibited as unsafe against it, for any
future change, which would make the specification worthless as a
regression net for exactly the class of bug it exists to catch.

\subsection{A decoupled demonstration}

To get a genuine counter-example, the safety property must not be
baked into the type the buggy operation is checked against. Two
small, self-contained Z documents were built and run independently
(neither is part of this file's own checked namespace, to avoid
exactly the trap above and to avoid redeclaring \field{Delegate} and
\field{Abandon} a second time in one namespace). Both share a
stripped-down state schema \field{MissionU} carrying only
\field{status}, \field{budgetRounds}, \field{blockingDelegationCount},
\field{resultRounds}, and \field{repoRootSet} --- structural bounds
only, no abandon invariant --- plus \field{Init} and \field{Delegate}
schemas identical in shape to \S\ref{sec:init}--\ref{sec:ops} (renamed
to \field{blockingDelegationCount} for consistency with the current
main document; the rename has no bearing on this demonstration, which
predates and is orthogonal to DES-076 --- see the note opening
\S\ref{sec:validation}). Each document then adds one more schema:

\begin{verbatim}
Idle
  Xi MissionU
where
  not (status = stAbandoned and
    (blockingDelegationCount /= 0 or resultRounds /= {}))
\end{verbatim}

\field{Idle} is a self-loop enabled everywhere \emph{except} the one
state this demonstration exists to find: abandoned with a live
blocking delegation or an unclaimed result. That converts the safety
property into a reachability question ProB already knows how to
answer without a custom goal: if the buggy operation can reach that
state, nothing is enabled there --- \emph{deadlock}.

The first document pairs \field{Init}/\field{Delegate}/\field{Idle}
with \field{AbandonFailOpen} (as given above, retyped against
\field{MissionU}). \texttt{/z-spec:check} passes;
\texttt{/z-spec:model\_check} explores the decoupled
\field{AbandonFailOpen} state space reported in the Tool run table and
reports a deadlock directly (not merely via a separate coverage or
liveness check), with this counter-example trace:

\begin{verbatim}
1  SETUP_CONSTANTS
2  INITIALISATION      (repoRootSet := zfalse)
3  Delegate            (blockingDelegationCount: 0 -> 1)
4  AbandonFailOpen      status: stOpen -> stAbandoned
                        (guard vacuous: repoRootSet = zfalse)
   -- deadlock: no operation enabled --
   status = stAbandoned, blockingDelegationCount = 1
\end{verbatim}

This is the fail-open bug, reached in four steps: delegate a worker,
then abandon under a \field{Store} with no \field{repoRoot} in
scope, and the mission is left \field{stAbandoned} with a live
blocking delegation record no longer reachable through any operation
in the model --- the formal mirror of the real bug silently discarding
recoverable work.

The second document is identical except \field{Abandon}
(\S\ref{sec:opabandon}'s guard, retyped against \field{MissionU})
replaces \field{AbandonFailOpen}. \texttt{/z-spec:check} passes;
\texttt{/z-spec:model\_check} explores the decoupled \field{Abandon}
state space reported in the Tool run table and reports \textbf{no
assertion or invariant violation} --- the state \field{Idle} excludes
is not reached in this exploration, because \field{Abandon}'s
$\field{repoRootSet} = \field{ztrue}$ conjunct is not an implication
and so is never vacuously satisfied. Unlike the buggy document's run,
this one is flagged ``not certified complete'' rather than passing
outright (\S\ref{sec:probrun}): probcli's bounded search does not
warrant that every valuation \field{MissionU}'s own type predicate
admits was actually enumerated, only that none it did reach violates
anything. The asymmetry is itself informative and is reported here
rather than smoothed over --- the buggy document's deadlock is a
positive, reproduced finding; the correct document's clean run is the
absence of one, at the bound probcli actually explored, not a
certificate of deadlock-freedom.

\subsection{Tool run}
\label{sec:probrun}

\begin{center}
\begin{tabularx}{\textwidth}{@{}lccl@{}}
\toprule
\textbf{Model} & \textbf{States} & \textbf{Transitions} &
  \textbf{Result} \\
\midrule
Main document (\S\ref{sec:state}--\ref{sec:ops}), all 9 operations
  $+$ \field{Idle} & 373 & 989 & 0 violations \\
Same, $+$ \field{AbandonFailOpen} (errata above) & 385 & 1025 &
  0 violations --- false negative, see errata \\
Decoupled, \field{AbandonFailOpen} & 7 & 21 & deadlock at
  $\field{status} = \field{stAbandoned} \land
  \field{blockingDelegationCount} = 1$; a deadlock search halts at the
  first violation, so this count is not the full space \\
Decoupled, \field{Abandon} & 15 & 29 & 0 violations, not certified
  complete \\
\bottomrule
\end{tabularx}
\end{center}

The two decoupled counts above are search-order artifacts, not
reproducible constants: probcli's deadlock search stops at the first
violation it finds, so the exact state/transition count depends on the
order it happens to explore the space in, which can vary across runs,
tool versions, or even independent re-implementations of the same two
schemas (\fpath{spec-writeset-admission.tex}'s own Tool run table
notes the identical caveat for its decoupled rows). What is
reproducible, and is the entire point of this demonstration, is the
\emph{verdict}: the buggy \field{AbandonFailOpen} pairing always
deadlocks, at the state
$\field{status} = \field{stAbandoned} \land
\field{blockingDelegationCount} = 1$
reached by the same four-step trace given below, regardless of exactly
how many states probcli visits before reporting it.

fuzz's \texttt{/z-spec:check} type-checks this document with 0 errors
in every revision recorded in its commit history. \texttt{probcli}
requires a \texttt{.tex} extension to run its Z-to-B translation ---
\texttt{/z-spec:check} runs fuzz directly on any extension, but
\texttt{/z-spec:test}/\texttt{/z-spec:model\_check} silently produced
0 states and 0 transitions when an earlier revision of this content
was saved with a \texttt{.md} extension --- the file was renamed to
\texttt{.tex} for exactly this reason before this document was first
submitted for review. All \texttt{model\_check}/\texttt{animate} runs
reported in this section, including the two decoupled documents, were
re-run against the current \texttt{.tex} file at this re-baseline.
The Result column reports the \texttt{model\_check} outcome
specifically (assertion/invariant violations), matching how this
section has always been read; \S\ref{sec:probrun-coverage} records a
separate, pre-existing finding the fuller \texttt{/z-spec:test}
checklist surfaces that \texttt{model\_check} alone does not.

\subsection{A disclosed, pre-existing gap: \field{AdvanceRound}
coverage}
\label{sec:probrun-coverage}

Running the fuller \texttt{/z-spec:test} checklist (init, animate,
assertions, deadlock-freedom, model-check, operation coverage) against
the main document --- both the current, re-baselined operation set
and, checked directly against the pre-re-baseline commit for
comparison, the original seven-operation set --- reports the same
coverage failure in both: \emph{\field{AdvanceRound} never fires}
within the explored 373-state space. This is not a defect introduced
by this re-baseline: it reproduces identically, state-for-state and
transition-for-transition (373 states, 749 transitions on the
original operation set), against the commit this document itself was
pinned to before today. Nothing in the operations this re-baseline
touches (\field{Abandon}, \field{DelegateAborted},
\field{DisclaimDelegation}) constrains \field{AdvanceRound}'s guard or
\field{budgetRounds}/\field{reflected} in any way, so there is no
DES-076-shaped fix available here; diagnosing and, if warranted,
fixing this is scoped as separate follow-up work, not folded into this
DES-076/DES-072 re-baseline, and is disclosed here rather than left
for the next reader to rediscover. \field{AdvanceRound} is not
otherwise unreachable by construction --- its guard
($\field{currentRound} < \field{budgetRounds} \land \field{currentRound}
\in \dom \field{reflected} \land \field{reflected(currentRound)} \notin
\field{TerminalRecommendation}$) is satisfiable on paper --- so the gap
most likely lies in how \field{Init} leaves \field{budgetRounds}$'$
unconstrained (\S\ref{sec:init}) interacting with probcli's
non-exhaustive bounded search, not in a logical impossibility; that
diagnosis is itself part of the deferred follow-up, not asserted here
as settled.

% ===================================================================
\section{Traceability}
\label{sec:trace}
% ===================================================================

\begin{center}
\begin{tabularx}{\textwidth}{@{}llX@{}}
\toprule
\textbf{Z construct} & \textbf{Go source} & \textbf{Note} \\
\midrule
\field{Status} & mission.go:10--28 & \field{validStatuses},
  validate.go:28--34 \\
\field{Recommendation} & reflection.go:89--103 &
  \field{validRecommendations} \\
\field{Mission} & \field{Contract} (mission.go:35--174) $+$
  blocking delegation count $+$ result/reflection round sets & the
  real \field{Contract} struct does not itself carry delegation or
  result data --- those live in sibling files (store.go:1959--1999,
  3450--3473) --- but the lifecycle invariants span all three, so the
  model's state includes them \\
\field{Init} & \field{Store.Create}
  (store.go:606--728) via \field{ApplyServerFields}
  (store.go:323--368) & \\
\field{Delegate} & \field{WriteDelegationSkeleton}
  (delegation.go:346--396) & status guard is a modelling choice, see
  \S\ref{sec:opdelegate} \\
\field{DelegateAborted} & \field{CloseDelegationSkeleton}
  (delegation.go:649--726) called from \field{closeDelegationAborted}
  (internal/hook/pretooluse\_dispatch.go:991--1014) or
  \field{closeSkeletonOnHashRefusal}
  (internal/hook/subagent\_start.go:787--833); read side
  \field{countBlockingDelegations} (store.go:1993--1995) & new,
  DES-076; \Xi Mission, see \S\ref{sec:opdelegateaborted} \\
\field{DisclaimDelegation} & \field{Store.DisclaimDelegation}
  (store.go:1516--1633) via \field{DisclaimDelegationRecord}
  (delegation.go:726--790) & new, DES-076; per-delegation eligibility
  not representable at this count-based granularity, see
  \S\ref{sec:opdisclaim} \\
\field{SubmitResult} & \field{Store.AppendResult}
  (store.go:3334--3450) & \\
\field{Reflect} & \field{Store.AppendReflection}
  (store.go:2703--2807) & \\
\field{AdvanceRound} & \field{Store.AdvanceRound}
  (store.go:2878--2963) & \\
\field{Close} & \field{Store.Close} (store.go:956--1187) via
  \field{checkResultGateLocked} (store.go:3516--3533) & \\
\field{Abandon} & \field{Store.Abandon} (store.go:1290--1487) &
  post-repoRoot-fix and post-DES-076; see abandon\_test.go:151--179
  and DESIGN.md's DES-076 entry \\
\field{AbandonFailOpen} & historical \field{Store.Abandon}, fixed
  PR \#432 (9a60b39) & not present in current code; modelled only
  for \S\ref{sec:validation} \\
(excluded) \field{Correct} & \field{Store.Correct}
  (correct.go:187--228) & DES-072; excluded by design, see
  \S\ref{sec:des072} \\
\bottomrule
\end{tabularx}
\end{center}

Pinned baseline commit: \textbf{244b6c3} (PR \#517), superseding the
original \textbf{9a60b39} (PR \#432) pin. See \S\ref{sec:currency} for
the currency rule this re-baseline exists to satisfy going forward.

\end{document}
